\documentclass[11pt,a4paper]{article}

\usepackage[a4paper,margin=1in]{geometry}
\usepackage[T1]{fontenc}
\usepackage[utf8]{inputenc}
\usepackage{lmodern}
\usepackage{xcolor}
\usepackage{hyperref}
\usepackage{enumitem}
\usepackage{tabularx}
\usepackage{array}

\hypersetup{
  colorlinks=true,
  linkcolor=blue!50!black,
  urlcolor=blue!50!black,
  pdftitle={A Text-Native Operating Filesystem for Natural-Language Robotic-Agentic Programming},
  pdfauthor={MD-OS (Artificial Prefrontal Cortex) v5.0},
  pdfsubject={Operational Context as Filesystem},
  pdfkeywords={MD-OS, Operating Filesystem, MCP, Codex, robotic-agentic programming, natural-language programming, persistent agents, AGI-like operation}
}

\setlist[itemize]{leftmargin=1.5em,itemsep=0.25em,topsep=0.25em}
\setlist[enumerate]{leftmargin=1.7em,itemsep=0.25em,topsep=0.25em}
\renewcommand{\arraystretch}{1.25}
\emergencystretch=3em

\newcommand{\code}[1]{\texttt{#1}}
\newcommand{\layerbox}[2]{%
  \fbox{\begin{minipage}{0.88\linewidth}
  \textbf{#1}\\[-0.1em]
  #2
  \end{minipage}}}

\title{\textbf{A Text-Native Operating Filesystem for Natural-Language Robotic-Agentic Programming}\\
\large Operational Context as Filesystem Through Markdown, Structured State, MCP, and Deterministic Builders}
\author{MD-OS (Artificial Prefrontal Cortex) v5.0}
\date{April 26, 2026}

\begin{document}
\maketitle

\begin{abstract}
This paper presents MD-OS (Artificial Prefrontal Cortex) v5.0, an early reference implementation of a
Markdown-native Operating Filesystem for natural-language robotic-agentic
programming. The central claim is that complex agentic and robotic ecosystems
should be programmed through durable, readable operating artifacts rather than
hidden chat history, opaque orchestration memory, or isolated tool calls.
Operational context is externalized into a filesystem composed of Markdown
knowledge, structured JSON state, append-only NDJSON logs, deterministic
builders, and bounded connectors. In this release, Codex is a verified
host runtime role and MCP is the tool and resource surface used to bind
the operating context to existing work systems, applications, devices, sensors,
and robots. MD-OS is not AGI and does not claim to create AGI; it supports
AGI-like operation in the practical sense of persistent memory, planning,
bounded tool use, audit, replay, role-bounded action, and human correction.
\end{abstract}

\section{Introduction}

Most agentic systems are still framed either as stateless chats with tool calls
or as hidden orchestration engines whose real state is difficult for a human to
inspect. Neither framing matches operational work. Real work is performed
through notes, procedures, tickets, documents, calendars, mail, planning boards,
filesystems, applications, terminals, APIs, devices, and human approvals.

MD-OS starts from a different premise:

\begin{quote}
The operational context of an agent should be externalized into a filesystem
that humans and machines can read, audit, reconstruct, transfer, and act from.
\end{quote}

This does not replace deterministic software with prose. It treats natural
language as a persistent control layer while deterministic modules perform
bounded state mutations, validation, replay, and connector execution.

\section{Core Thesis}

MD-OS (Artificial Prefrontal Cortex) v5.0 is a text-native, auditable, and reconstructible control
plane for persistent agents and robotic systems. It is best described as an
Operating Filesystem: a filesystem-resident operating context that surrounds a
host model runtime.

The core thesis is:

\begin{quote}
AI memory is not chat history. AI action is not raw tool calling. Persistent
agent operation needs an operating filesystem: durable knowledge, explicit
state, bounded tools, replay, audit, and human correction.
\end{quote}

MD-OS is not a Linux replacement, not ROS, not firmware, not a real-time
hardware control loop, and not a high-frequency OLTP database. It is the
filesystem-native control plane around host reasoning and bounded tool use.

\section{Natural-Language Robotic-Agentic Programming}

The most important paradigm is not simply file-based memory. It is
natural-language programming of a complex robotic-agentic ecosystem.

The target system is not a single chatbot or one robot command. The target is
the ecosystem around work:

\begin{itemize}
  \item humans, experts, operators, and approval boundaries;
  \item host runtimes such as Codex;
  \item MCP resources and tools;
  \item mail, calendar, agenda, tickets, documents, planning boards, and
  internal applications;
  \item terminals, APIs, queues, services, filesystems, and databases;
  \item devices, sensors, robot controllers, actuators, and telemetry streams;
  \item policies, procedures, exceptions, recovery paths, and audit records.
\end{itemize}

MD-OS programs that ecosystem through durable operating artifacts:

\begin{quote}
mission / role / policy / procedure / exception
$\rightarrow$ Markdown operating artifact
$\rightarrow$ deterministic builder or compiler
$\rightarrow$ bounded connector action or observation
$\rightarrow$ telemetry, snapshot, or audit event
$\rightarrow$ updated operating memory
\end{quote}

This is why natural language matters. It is not treated as transient prompt
text. It becomes the readable programming surface for missions, roles,
procedures, safety boundaries, connector contracts, permissions, escalation
rules, evaluations, and recovery procedures.

\section{Layer Stack}

The current 5.0 release is organized as a layered system.

\begin{center}
\begin{tabular}{>{\centering\arraybackslash}p{0.92\linewidth}}
\layerbox{Human Supervision Layer}{new hire, expert, operator, or supervisor;
intent, approval, correction, accountability}\\[0.5em]
$\downarrow$ intent / chat\\[0.3em]
\layerbox{Codex Host Runtime}{execution host; LLM chat, tool
orchestration, bootstrap readback; execution layer, not canonical memory}\\[0.5em]
$\downarrow$ state readback\\[0.3em]
\layerbox{MD-OS Operating Filesystem}{active boundary: \code{md-os/}; persistent
identity, role context, state, agenda, audit, replay; Markdown, JSON, and NDJSON
as readable operating truth}\\[0.5em]
$\downarrow$ compile / rebuild\\[0.3em]
\layerbox{Deterministic Runtime}{\code{md-os/os/*.js}; builders, validators,
compilers, role intake, role sensemaking, replay, indexes}\\[0.5em]
$\downarrow$ bounded resources and tools\\[0.3em]
\layerbox{MCP and Connector Boundary}{resources, tools, snapshots, permission
profiles, bounded reads and writes to authorized work sessions}\\[0.5em]
$\downarrow$ external substrates\\[0.3em]
\layerbox{Host OS, Applications, Work Systems, APIs, Devices}{Linux, macOS,
Windows, shell, files, desktop applications, browsers, mail, calendar, agenda,
planning boards, tickets, documents, internal applications, queues, services,
sensors, robots, and telemetry}
\end{tabular}
\end{center}

Codex can supply reasoning in a verified interactive workflow. MD-OS supplies durable
operating context. MCP supplies the resource and tool surface through which the
host can observe or act on existing substrates.

\section{Natural Language as Program}

Natural language becomes operationally serious when it defines:

\begin{itemize}
  \item what the system is allowed to do;
  \item what order of operations is valid;
  \item how entities and states should be named;
  \item what depends on what;
  \item which exceptions are known;
  \item what requires approval;
  \item and what counts as completion.
\end{itemize}

Markdown is useful because it can encode procedures, guardrails, decision
criteria, role contracts, connector contracts, boundary conditions, and
cross-references between operational concepts. It behaves as a high-level
batch language for agentic work. The Markdown files do not execute like shell
scripts; they define operating semantics that deterministic builders compile
or materialize into structured runtime state.

\section{Operating Filesystem Structure}

The active operational boundary is \code{md-os/}. The important sublayers are:

\begin{itemize}
  \item \code{md-os/kb/}: stable knowledge, architecture, permission models,
  connector contracts, runtime lifecycle models, and cognitive bootstrap;
  \item \code{md-os/ops/}: persistent runtime state, role intake, project state,
  agendas, summaries, journal, connector registry, and generated views;
  \item \code{md-os/os/}: deterministic scripts, builders, validators, compilers,
  replay, connector executors, and MCP server adapter;
  \item \code{md-os/schemas/}: JSON contracts for structured runtime state.
\end{itemize}

This makes the system inspectable. A human can review the same state that the
host runtime reads. A later session can reconstruct orientation by reading
files rather than trusting hidden process memory.

\section{Codex as a Verified Host Path}

MD-OS is host-portable in architecture. Codex is one verified execution layer
through which MD-OS can be read, operated, modified, and reported in an
interactive workflow.

This distinction matters:

\begin{itemize}
  \item MD-OS is the persistent agent identity and operating context on disk.
  \item Codex is the current reasoning and execution host.
  \item MCP is the tool and resource surface.
  \item The host can change only after bootstrap, permissions, command
  behavior, working directory behavior, and runtime readback are verified.
\end{itemize}

The cognitive bootstrap is therefore not branding. It is the procedure that
frames the host runtime so it answers from the MD-OS operating context rather
than from a generic stateless chat session.

\section{MCP, Connectors, and Internal Work Systems}

MD-OS is not API-first. It is substrate-first. APIs are valuable when they
exist, but many organizations operate through mail, calendars, spreadsheets,
desktop applications, browsers, shared folders, terminals, ticketing systems,
planning boards, queues, and legacy tools.

MCP and connectors let MD-OS bind to those existing resources without requiring
a bespoke API project before useful onboarding or assistance can begin. In a
workplace deployment, the new hire can use already-authenticated and
authorized sessions for:

\begin{itemize}
  \item mail and calendar;
  \item agenda and planning tools;
  \item tickets, requests, and queues;
  \item documents and shared folders;
  \item internal applications;
  \item terminals, APIs, and databases where allowed;
  \item devices, robots, or telemetry surfaces under explicit safety policy.
\end{itemize}

The authority boundary stays with the human and the host environment. MD-OS
does not claim hidden credentials or unsupervised internal authority. It
records context, permissions, evidence, snapshots, and next actions inside the
operating filesystem.

\section{Role Onboarding and Operational Sensemaking}

A major practical use case is new-hire assistance. Companies often cannot hand
over a clean training corpus. They hand over a pile of files, exports, notes,
PDFs, spreadsheets, old procedures, emails, screenshots, and examples of past
work.

MD-OS models this through role-first intake:

\begin{itemize}
  \item \code{md-os/ops/roles/<role\_id>/ROLE.md}: role contract;
  \item \code{md-os/ops/roles/<role\_id>/intake/raw/}: messy raw material;
  \item \code{cortex role intake <role\_id>}: inventory and weak extraction;
  \item \code{cortex role sensemake <role\_id>}: reconstructed cases, relation
  graph, root-cause candidates, work patterns, and expert questions.
\end{itemize}

The goal is not to pretend that the system instantly knows the company. The
goal is to reduce training time, repeated internal errors, dependency on tacit
knowledge, and internal integration cost. The new hire works in Codex chat,
with MD-OS APFC as a role assistant grounded in the role contract, reconstructed
cases, work patterns, and MCP-visible work tools.

\section{AGI-like Operation Without an AGI Claim}

MD-OS should not be described as AGI. It does not claim consciousness,
sentience, or general intelligence. The intelligence comes from the host model,
the human-supervised loop, and the operating artifacts.

It is accurate to say that MD-OS supports AGI-like operation in a practical
engineering sense:

\begin{itemize}
  \item persistent memory across sessions;
  \item explicit goals and roles;
  \item bounded tool use;
  \item planning and task state;
  \item replayable audit;
  \item recovery after interruption;
  \item correction through normal files;
  \item and progressive evolution through knowledge, programs, builders,
  schemas, tests, and connectors.
\end{itemize}

This is an operating-context claim, not a claim that MD-OS itself is generally
intelligent.

\section{Replay and Deterministic Builders}

The runtime is designed around deterministic builders. The host runtime may
interpret intent, propose changes, or call tools, but canonical state is
materialized by scripts under \code{md-os/os/}.

The replay pattern is central:

\begin{quote}
source signals + programs + journal + deterministic builders $\rightarrow$
compiled state + agendas + summaries + indexes
\end{quote}

This makes continuity testable. If generated views can be removed and rebuilt
from sources, the system does not rely on trust in the previous chat session.

\section{Production Limits}

MD-OS is strongest today in supervised, auditable, role-bounded workflows. It
is not yet unsupervised production infrastructure. The main production limits
are:

\begin{enumerate}
  \item \textbf{I/O scalability.} Text files are excellent canonical truth, but
  high-volume query paths need derived indexes, sharding, summaries, and
  compaction.
  \item \textbf{Security boundaries.} Markdown policy is auditable governance,
  not hard security. Hard enforcement must live in host permissions, sandboxing,
  MCP broker rules, read-only files, allowlists, and external infrastructure.
  \item \textbf{Host dependency.} MD-OS externalizes memory but still depends on
  a reasoning runtime. Host compatibility must be verified per host path.
  \item \textbf{UX complexity.} The underlying model is rigorous. The user
  surface must expose high-level commands rather than forcing humans to manage
  every registry and schema manually.
  \item \textbf{Multi-agent concurrency.} Local file locks and append-only
  proposals are enough for supervised workflows, but fleet-scale multi-agent
  operation requires queues, leases, ownership, and external conflict
  coordination.
\end{enumerate}

\section{Conclusion}

The contribution of MD-OS is not a claim of full autonomy. It is a
natural-language programming substrate for complex agentic and robotic
ecosystems, made operationally serious through text-native structure.

Natural language becomes a practical programming and control medium when it is:

\begin{itemize}
  \item stabilized into auditable artifacts;
  \item bounded by deterministic executors;
  \item connected to persistent state;
  \item exposed to existing work systems through MCP and connectors;
  \item and constrained by explicit human-supervised operating rules.
\end{itemize}

Under those conditions, Markdown, JSON, NDJSON, deterministic scripts, MCP
resources, and connector snapshots together form a practical Operating
Filesystem for programming persistent AI agents, workplace role assistance,
robotic systems, devices, and host runtimes as one supervised ecosystem.

\end{document}
